\section{Notation \tbc}

Although the entire construction described in the writeup applies to arbitrary prime fields, we confine ourselve to the the following choices.
%
Let $\mathbb F_q$ be the Goldilocks prime field \cite{}, with the Solinas prime
\[
q = 2^{64} - 2^{32} + 1
\]
as modulus. 
We denote by $F\geq \mathbb F_q$ any cryptographically large extension, in our case the degree 3 extension $F = \mathbb F_{q^5}$, using...
\ulrich{%
Check which irreducible polynomial is best
}


% \subsection{Notation}
Boolean Hypercube $H_n = \{0,1\}^n\subset \mathbb F_q^n$, indexing in natural order via $[0,2^n)$, 
multivariate and multilinear polynomials,
Lagrange polynomial
\begin{equation}
\label{e:eq}
eq(\vec y, \vec\rho) = \prod_{i=1}^n (y_i \cdot \rho_i + (1 - y_i)\cdot (1-\rho_i))
%= \prod_{i=1}^m (1 - y_i - \rho_i + 2 y_i  \rho_i),
\end{equation}

Inner product notation.


\subsection{Multilinears and  univariates}

 
Given a function $f\in \mathbb F^{H_n}$ (a vector of values) we will write
\[
f(X_1,\ldots, X_n) \in \mathbb F[X_1,\ldots, X_n]^{<2}
\]
for its multilinear extension, and 
% multilinear $f(X_1,\ldots, X_n)\in $
% %, given by its values over the Boolean hypercube $H_n$, %(the component of an interleaved representation) with the univariate polynomial
associate to it the univariate polynomial
\begin{equation}
\label{e:values:to:coeffs}
u_f(X) = \sum_{i=0}^{2^n-1} f(i_1,\ldots, i_n) \cdot X^i,
\end{equation}
defined by the values of $f$ over $H_n$,  where $i = i_1 + \ldots + 2^{n-1} i_n$ is the decomposition of the index into bits. 
Under this \textit{values-to-coefficient} regime, we identify the space of multilinears (in $n$ variables) with the univariate space
\begin{equation}
\mathscr P_{n}(\mathbb F) = \mathbb F[X]^{<2^n} 
\end{equation}
in a way that is particularly efficient for Reed--Solomon encoding. 

% \subsubsection{Liftings.}
% Given a multilinear $f$ in $k$ variables, and $n>k$, we call
% \[
% f^*(X_1,\ldots, X_n) = (1-X_1)\cdot\ldots (1-X_{n-k})\cdot f(X_{n-k+1},\ldots, X_n)
% \]
% its \textit{lifting} to $n$ variables. 
% Lifting will allow to treat a collection of polynomials in a uniform manner, while not paying the cost of duplicated values.
% In terms of univariate representations,
% \[
% u_{f^*}(X) = u_f(X^{2^{n-k}}).
% \]

\subsubsection{Univariate evaluation}
Evaluation at a point $z\in\mathbb F$ is a linear functional and can be written as inner product $u_m(z) = \langle f, L_z\rangle_{H_n}$ with  
\begin{equation}
\label{e:uni:eval:kernel}
L_z(X_1, \ldots, X_n) = \prod_{j=1}^n \left(
 (1 - X_j) + X_j \cdot z^{2^{j-1}}  \right),
\end{equation}
and its values over the hypercube are the powers of $z$ in natural order.
% In terms of codes, $L_z$ expresses point evaluation of the codeword as linear functional on the message.

% This establishes a one-to-one correspondence between multilinear polynomials in $k$ variables, and univariate polynomials from the space 
% \[
% \mathscr P_{k} = \mathbb F[X]^{<2^k}.
% \]

% For example the univariate Lagrangian 
% \[
% L_z(X) = \frac{1}{2^n}\cdot \left(1 + \sum_{i=1}^{2^n-1} z^{2^n - i} \cdot X^i\right) 
% \]

% Interleaved representations are encoded component-wise, i.e. row by row, using a Reed--Solomon code with message length matching the row length. 
% Concretely, each postfix function $g(X_1, \ldots, X_{n-k}) = f(X_1,\ldots, X_{n-k},\vec y)$, given by its values over the Boolean hypercube $H_{n-k}$, is associated with the univariate polynomial
% \[
% \sum_{\vec x\in \{0,1\}^{n-k}} g(\vec x) \cdot X^{\iota(\vec x)} \quad \in \mathbb F[X]^{<2^{n-k}},
% \]
% where $\iota(\vec x) = x_1 + 2 x_2 + \ldots + 2^{n-k-1} x_{n-k}$ is the natural index of $\vec x$, and evaluated over a specified domain $D\subset \mathbb F$. 
% In this way, we obtain postfix many Reed--Solomon codes, viewed as word from the \textit{interleaved code}.






\subsection{Interleaved representation}

In alignment with the sumcheck protocol, we will often think of a multilinear $f(X_1,\ldots, X_n)$ decomposed into its $k$-prefix functions
% by postfixes $\vec y\in \{0,1\}^k$,
\begin{align}
\label{e:prefix:decomp}
f(X_1,\ldots,X_n) =  \sum_{i = 0}^{2^{k}-1} eq_{i}(X_{1},\ldots, X_k) \cdot f_i(X_{k+1},\ldots, X_{n}) \cdot 
\end{align}
Here we again used the identification  $H_k = [0,2^k)$, by each index $i$ specifying the prefix $(i_{1},\ldots, i_{k})\in H_{k}$ via its bits.
We call 
$
(f_0, \ldots, f_{2^{n-k}-1}) %\in \mathbb F[X_{k+1},\ldots, X_n]^{2^{k}}
$
the \textit{interleaved representation} of $f$, with interleaving depth $e = 2^{k}$.
In terms of their univariate polynomials,
\begin{align}
\label{e:FFT:decomp}
u_f(X) = \sum_{i=0}^{2^k-1}  X^{i} \cdot u_{f_i}(X^{2^k}), 
\end{align}
which is the FFT decomposition of $u_f(X)$. 

\begin{rem}%[Postfix decomposition]
\label{rem:postfix:decomp}
The current version of ProveKit instead employs a postfix decomposition of multilinear polynomials, which corresponds, in the univariate representation, to partitioning the polynomial into contiguous blocks of monomial powers.
%$u_f(X) = \sum_{i=0}^{2^{n-k}-1} u_{f_i}(X) \cdot X^{i \cdot 2^{k}}$.
Since our construction relies very little on properties specific to Reed--Solomon codes, the choice of decomposition is largely a matter of convenience. We adopt a prefix decomposition instead, as it leads to notation that is more consistent with our treatment of liftings.
\end{rem}



\subsection{Liftings}

Given a multilinear $f$ in $k$ variables, and $n>k$, we call
\begin{equation}
\label{e:lifting:multilinear}
f^*(X_1,\ldots, X_n) = (1 - X_1)\cdot \ldots \cdot (1 - X_{n-k}) \cdot f(X_{n-k+1},\ldots, X_n)
\end{equation}
its \textit{lifting} to $n$ variables. 
In terms of the values over the hypercube, lifting simply places the original vector (the values over $H_k$) on the prefix cube 
\[
\{0^{n-k}\}\times H_{k} \subset H_n,
\]
and pads it with zero on the remainder of $H_n$.
In terms univariate representations, 
\begin{equation}
u_{f^*}(X) = u_f(X^{2^{n-k}}).
\end{equation}
We use lifting to uniformize our R1CS across the two stages at no additional cost.


% \subsection{Sumcheck}

% % In reverse order, since the hypercube is indexed in natural order.

% The multivariate sumcheck protocol \cite{Sumcheck} is an interactive proof for that a multivariate polynomial $G\in \mathbb F[X_1,\ldots, X_n]$ satisfies
% \[
% \sum_{\vec x \in H_n} G(\vec x) = s.
% \]
% % Typically, the multivariate $G$ is a virtual polynomial in several committed multilinears, and its maximum individual degree $d = \max_i \deg_{X_i}(G)$ is at least $2$.
% We assume that the reader is familiar with it, and only point out the following:
% For implementation reasons, we consider the sumcheck protocol in \textit{reverse order}, from the last variable $X_n$ down to the first variable $X_1$.
% In each step,  $k = n, n-1, \ldots, 1$, the prover provides the refinement polynomials $q_k(X)\in \mathbb F[X]^{\leq d_k}$, where $d_k$ is the individual degree of $G$ in $X_k$, and the verifier chooses to continue on the specialized claim $q_k(\lambda_k)$ at a random  $\lambda_k\sample\mathbb F$. 
% \[
% q_k(X) = \sum_{(x_{1},\ldots, x_{k-1}) \in H_{k-1}} G(x_1,\ldots, x_{k-1},X,\lambda_{k+1},\ldots, \lambda_n),
% \]


% Although we shall describe our protocols in a self-contained manner, we quickly sketch the sumcheck protocol, as it is a core ingredient of multilinear proof system.

% Before the first round the prover claims the refinement polynomial 
% \[
% q_n(X) = \sum_{(x_2,\ldots, x_{n-1}) \in H_{n-1}} G(x_1,\ldots, x_{n-1}, X),
% \]
% which is of degree at most $d$.
% The verifier checks the sanity condition that $s = q_0(0) + q_0(1)$, chooses a random $\lambda_n$, and it asks the prover to prove the specialized claim
% \[
% q_n(\lambda_n) = \sum_{(x_2,\ldots, x_n) \in H_{n-1}} G(x_1,\ldots, x_{n-1}, \lambda_n)
% \]
% instead.
% The protocol continues in the same manner, letting the prover provide refinement polynomials
% \[
% q_k(X) = \sum_{(x_{1},\ldots, x_{k-1}) \in H_{k-1}} G(x_1,\ldots, x_{k-1},X,\lambda_{k+1},\ldots, \lambda_n),
% \]
% for $i=1,\ldots,n-2$ as response of the received challenges $\lambda_2,\ldots, \lambda_{n-1}$.
% In the last step the verifier samples $\lambda_{n}\sample F$, which eventually reduces the sumcheck claim to that
% \[
% q_{n-1}(\lambda_n) = G(\lambda_1,\ldots, \lambda_n), 
% \]
% a evaluation claim for the multivariate polynomial $G$ at the random point $\vec \lambda = (\lambda_1,\ldots, \lambda_n)$, comprised of the random challenges received in the course of the protocol.
% For a detailed treatment of the sumcheck protocol, and its importance to multivariate proofs in general, we refer to the manuscript \cite{Thaler:book} and the references therein.



\subsection{Properties of interleaved Reed--Solomon codes}

Are RS codes over an extension field. 
Same rate, same distance.  

